% ECONOMICS_PROBLEM: {"id": "D24-460", "title": "Misallocation versus measurement error", "jel_code": "D24", "source_id": "OP-0262"}
% !TeX program = xelatex
\documentclass[11pt,a4paper]{article}
\usepackage[margin=23mm]{geometry}
\usepackage{fontspec}
\setmainfont{Latin Modern Roman}
\usepackage{amsmath,amssymb,mathtools}
\usepackage{xcolor,enumitem,booktabs,longtable,array,xurl,hyperref}
\definecolor{muted}{HTML}{53636C}
\hypersetup{colorlinks=true,urlcolor=blue,linkcolor=blue}
\setlength{\parindent}{0pt}
\setlength{\parskip}{5pt}
\newcommand{\R}{\mathbb R}
\newcommand{\E}{\mathbb E}
\newcommand{\N}{\mathbb N}
\newcommand{\ind}{\mathbf 1}
\newcommand{\Var}{\operatorname{Var}}
\newcommand{\Cov}{\operatorname{Cov}}
\newcommand{\supp}{\operatorname{supp}}
\DeclareMathOperator*{\argmin}{arg\,min}
\DeclareMathOperator*{\argmax}{arg\,max}
\newcommand{\entrymeta}[3]{{\sffamily\small\color{muted}\textbf{\texttt{#1}}\quad|\quad #2\par #3\par}}
\newcommand{\Problemref}[1]{\texttt{#1}}
\newcommand{\JELref}[2]{\texttt{#2} (JEL \texttt{#1})}
\begin{document}
\section*{Misallocation versus measurement error}
\textbf{Problem ID:} \texttt{D24-460}\quad \textbf{JEL:} \texttt{D24}\par
% BEGIN ECONOMICS STATEMENT
\label{problem:OP-0262}
{\sffamily\small\color{muted}Original subject group: Growth and productivity\par}
\label{definitions:problem:30007543}
\textbf{Audit classification:} Background; proposed extension. The source provides a measurement-error correction, not a declaration that all correction/identification is unsolved. Verification concerns the cited version, not first priority or proof that this exact editorial specification remains unsolved on October 8, 2026.
\subsection*{Definitions and precise research target}
\paragraph{Editorial mathematical specification.} Consider plants \(i\) producing a differentiated product with technology \(q_i=A_i k_i^\alpha \ell_i^\gamma m_i^{1-\alpha-\gamma}\), where capital, labor, and intermediates are positive, \(\alpha,\gamma>0\), \(\alpha+\gamma<1\), and production elasticities are specified or identified. Product demand has a stated elasticity, and observed revenue and inputs equal true values plus measurement errors in logs. Economic wedges may affect financing, labor, or intermediate purchases; adjustment costs and transport costs are modeled separately rather than automatically called distortions.
Require initial aggregate output \(Y>0\). Let \(Y^{\mathrm{eff}}\) be aggregate output after eliminating the specified wedges while holding true technologies, aggregate resources, and unavoidable costs fixed, and let \(Y\) be initial output. Define
\[
G=\frac{Y^{\mathrm{eff}}-Y}{Y}.
\]
Identify or bound \(G\) using repeated plant observations, independent audits, or instruments that distinguish measurement error from genuine marginal-product gaps. State error persistence, correlation, and demand assumptions explicitly; without them, dispersed measured revenue products need not identify misallocation. Quantify how much of observed dispersion comes from true wedges, measurement, demand heterogeneity, and adjustment costs. Validate the decomposition using independent responses to resource reallocations. The target is a defensible counterfactual productivity gain in a declared industry, not a presumption that all within-industry productivity dispersion indicates inefficient allocation.

Fix a final-good aggregator \(Y=(\sum_i\omega_iq_i^{(\epsilon-1)/\epsilon})^{\epsilon/(\epsilon-1)}\), \(\epsilon>1\), \(\omega_i>0\), together with aggregate capital, labor, and intermediate resources and a specified input-supply technology. The efficient benchmark maximizes this same quantity index over feasible allocations using true technologies, with only the declared wedges removed. Measurement errors satisfy \(\log\widetilde q_i=\log q_i+v_i^q\) and analogous input equations; specify their joint serial law rather than presuming independence. Under imperfect competition observed revenue products differ from true physical marginal products even with perfect measurement, and \(Y^{\rm eff}\) must state whether markups are removed too. The set for \(G\) consists of its values over all true-input, demand, technology, and error laws compatible with observed panels and independent validation restrictions. A measurement correction in one model does not settle identification under unrestricted error.
\subsection*{Variants and assumption changes}
The following are \emph{editorial variants}, not additional published open conjectures.
\begin{enumerate}
\item \emph{Independent validation.} Add repeated independent administrative measurements conditional on true plant inputs and output, plus a fixed CES demand system. Identify error variances and compare reallocation gains with the single-measurement case.
\item \emph{Persistent correlated error.} Allow measurement error to persist and covary with true productivity or plant size under quantified covariance bounds. Determine how wide the feasible gain set becomes and which extra audits shrink it.
\end{enumerate}
\subsection*{References and source audit}
\begin{itemize}
\item Mark Bils; Peter J. Klenow; Cian Ruane. \emph{Misallocation or Mismeasurement?}. 2020. Locator: Census CES20-07, February 2020, Abstract (method and results). Primary record and abstract/summary checked; no fulltext claim is made. \url{https://www.census.gov/library/working-papers/2020/adrm/CES-WP-20-07.html}.
\end{itemize}
The source provides a measurement-error correction, not a declaration that all correction/identification is unsolved.
\begin{enumerate}\raggedright
\item\hypertarget{definitions:source:30007543:1}{} Mark Bils; Peter J. Klenow; Cian Ruane. 2020. \emph{Misallocation or Mismeasurement?}. Census CES20-07, February2020, Abstract (method and results).
\url{https://www.census.gov/library/working-papers/2020/adrm/CES-WP-20-07.html}.
DOI: \url{https://doi.org/10.3386/w26711}.
\end{enumerate}

\subsection*{Common conventions: Source mathematical conventions}

These conventions apply to each entry unless explicitly replaced.

\textbf{Models and identification.} A primitive model \(m\) specifies all probability laws, feasible choices, information, equilibrium and measurement rules used by its target. The admissible class \(\mathcal M\) is fixed independently of outcomes. If \(P_O(m)\) is its observable probability law and \(\tau(m)\) the target, its identified set at observable law \(P\) is
\[
 \mathcal I_\tau(P)=\{\tau(m):m\in\mathcal M,\ P_O(m)=P\}.
\]
Point identification means that this set is a singleton. An empty set signals incompatibility of data and assumptions. Coordinate bounds need not describe the joint set. Bounds are sharp only if every retained target is attainable or arbitrarily approachable by compatible models. Finite bounds require explicit restrictions on unobserved quantities; unrestricted confounding, hidden wealth, extrapolation or arbitrary equilibrium selection can yield uninformative sets.

\textbf{Causal interventions.} A potential outcome \(Y_i(p)\) is the outcome under a complete policy regime \(p\), including timing, financing, information and market spillovers. The target population and units are fixed before treatment. With interference, use the complete assignment vector or a justified exposure mapping; individual no-interference notation alone cannot identify an economy-wide intervention. A difference of mean potential outcomes does not require the cross-world joint distribution. Conditioning on post-treatment migration, survival, enrollment or employment changes the estimand and needs separate justification.

\textbf{Statistical statements.} An estimator, confidence set, forecast or test specifies a sampling law, model class and loss/coverage criterion. Uniform \((1-\alpha)\)-coverage means \(\inf_{m\in\mathcal M}\Pr_m\{\tau(m)\in C_n\}\geq1-\alpha\), or an explicitly asymptotic version. The whole parameter space trivially has coverage, so informativeness requires a stated width, risk or power objective. A forecasting improvement is relative to a named benchmark, information set and evaluation population.

\textbf{Equilibrium and optimization.} An equilibrium comprises feasible plans, optimality, beliefs where needed, budget constraints and all market/resource clearing. The primitives or a stated benchmark define each correspondence \(\mathcal E(p;m)\). Nonemptiness is assumed or proved, never inferred just from compact policy sets. A selected-equilibrium target uses a declared selection rule; a robust target quantifies over all permitted equilibria. Use suprema or infima unless attainment is established. An \(\varepsilon\)-optimal policy is feasible and within \(\varepsilon\) of the finite optimum value. Report an empty feasible set separately.

\textbf{Welfare and accounting.} Normative utility, interpersonal normalization, social weights, numeraire, discounting and the use of public receipts are primitives. Behavioral utility need not coincide with the welfare criterion. Household welfare includes ownership income and rents once; adding profits again to compensating variation generally double-counts them. A monetary equivalent is defined by an explicit transfer and price convention, with monotonicity and range conditions. Average effects, mechanism contributions, transfers and welfare losses are different targets. Infinite sums require convergence and dynamic feasible plans require terminal or no-Ponzi conditions.

\textbf{Source versions.} A working-paper date, revision date and journal publication year are distinguished. A DOI for a journal version is not automatically the DOI of an earlier manuscript. Locators refer to the stated version and printed pages where specified. A metadata-only check does not verify a claim about a full-text conclusion. Every entry's source subsection narrows the provenance claim accordingly.
% END ECONOMICS STATEMENT
\end{document}
